% Chương 3 — phần bổ sung của tiểu luận final: CNN cài đặt bằng ba cách, CNN dạng hợp các hàm, đối chiếu ảnh tự nhiên.

\section{CNN cài đặt bằng ba cách}
\label{sec:ch3-three}

Các mô hình ở mục trước đều dùng PyTorch. Phần này viết \emph{cùng một} CNN cơ bản bằng ba cách: tự viết bằng NumPy, Keras (TensorFlow) và PyTorch. Mục đích là hiểu tích chập ở mức phép toán và kiểm chứng rằng ba bản cho cùng chất lượng. Khung huấn luyện kế thừa mã nguồn của thành viên nhóm 07, được mở rộng thêm bộ ảnh đáy mắt. Cả ba bản đọc cùng một tệp \texttt{retina.npz} (1.200 ảnh train+val gốc của RetinaMNIST, 400 ảnh test) và cùng một phép chia validation phân tầng 10\%, nên chênh lệch kết quả không đến từ dữ liệu.

\textbf{Kiến trúc cơ bản} (ảnh thu về 48×48×3): \texttt{Conv(3→32, 3×3) – ReLU – MaxPool(2) – Conv(32→64, 3×3) – ReLU – MaxPool(2) – Flatten – Dense(128) – ReLU – Dense(5) – softmax}, gồm 1.199.813 tham số (phần lớn nằm ở lớp Dense đầu tiên: $64\cdot12\cdot12\cdot128$). Ngoài ra có một \textbf{CNN cải tiến} (ảnh 96×96) cho hai framework có tự động đạo hàm: stem Conv–BN–ReLU, bốn khối dư (residual) 32–64–128–256 kênh xen MaxPool và Dropout tăng dần, Global Average Pooling, Dense(256), cùng tăng cường dữ liệu (lật, xoay ±15°, dịch, phóng, đổi màu).

\textbf{Bản NumPy} không dùng bốn vòng lặp lồng nhau mà dùng kỹ thuật \emph{im2col}: duỗi mỗi cửa sổ 3×3×C thành một hàng, khi đó tích chập trở thành \emph{một phép nhân ma trận} lớn (\coderef{lst:ch3-im2col}). Lan truyền ngược dùng phép ngược \emph{col2im} để cộng dồn gradient của các cửa sổ chồng lấn về đúng điểm ảnh.

\begin{code}{Tích chập bằng im2col và lan truyền ngược (\texttt{tl3/ch3\_scratch.py}, rút gọn)}
\label{lst:ch3-im2col}
\begin{minted}{python}
def im2col(X, kh, kw, stride=1, pad=0):
    """X [N,C,H,W] -> [N·Ho·Wo, C·kh·kw]: mỗi hàng là một cửa sổ đã duỗi thẳng; chỉ lặp kh·kw lần."""
    N, C, H, W = X.shape
    Ho, Wo = (H + 2 * pad - kh) // stride + 1, (W + 2 * pad - kw) // stride + 1
    Xp = np.pad(X, ((0, 0), (0, 0), (pad, pad), (pad, pad))) if pad else X
    cols = np.empty((N, C, kh, kw, Ho, Wo), dtype=X.dtype)
    for i in range(kh):
        for j in range(kw):
            cols[:, :, i, j] = Xp[:, :, i:i + stride * Ho:stride, j:j + stride * Wo:stride]
    return cols.transpose(0, 4, 5, 1, 2, 3).reshape(N * Ho * Wo, -1), Ho, Wo

def conv_forward(X, W, b, pad=1):
    F = W.shape[0]
    cols, Ho, Wo = im2col(X, W.shape[2], W.shape[3], 1, pad)
    Z = cols @ W.reshape(F, -1).T + b                        # tích chập = một phép nhân ma trận
    return Z.reshape(X.shape[0], Ho, Wo, F).transpose(0, 3, 1, 2), (X.shape, cols, W, pad, Ho, Wo)

def conv_backward(dZ, cache):
    X_shape, cols, W, pad, Ho, Wo = cache
    dZ2 = dZ.transpose(0, 2, 3, 1).reshape(-1, W.shape[0])
    dW = (dZ2.T @ cols).reshape(W.shape)                     # dL/dW = dZ^T · cols
    db = dZ2.sum(axis=0)
    dX = col2im(dZ2 @ W.reshape(W.shape[0], -1), X_shape, *W.shape[2:], Ho, Wo, 1, pad)
    return dX, dW, db
\end{minted}
\end{code}

\begin{code}{Cùng CNN cơ bản bằng Keras (\texttt{tl3/ch3\_keras.py}) và PyTorch (\texttt{tl3/ch3\_torch.py})}
\label{lst:ch3-keras-torch}
\begin{minted}{python}
keras_cnn = keras.Sequential([layers.Input((48, 48, 3)),
    layers.Conv2D(32, 3, padding="same", activation="relu"), layers.MaxPooling2D(2),
    layers.Conv2D(64, 3, padding="same", activation="relu"), layers.MaxPooling2D(2),
    layers.Flatten(), layers.Dense(128, activation="relu"), layers.Dense(5, activation="softmax")])

torch_cnn = nn.Sequential(
    nn.Conv2d(3, 32, 3, padding=1), nn.ReLU(), nn.MaxPool2d(2),
    nn.Conv2d(32, 64, 3, padding=1), nn.ReLU(), nn.MaxPool2d(2),
    nn.Flatten(), nn.Linear(64 * 12 * 12, 128), nn.ReLU(), nn.Linear(128, 5))   # softmax trong CrossEntropyLoss
\end{minted}
\end{code}

\begin{table}[H]
\centering
\caption{CNN ba cách trên RetinaMNIST (400 ảnh test; luôn đoán mức 0 đạt accuracy 0,435)}
\label{tab:ch3-three}
\small
\begin{tabular}{@{}llcccrr@{}}
\toprule
\textbf{Cách} & \textbf{Mô hình} & \textbf{Accuracy} & \textbf{Macro-F1} & \textbf{QWK} & \textbf{Tham số} & \textbf{Huấn luyện} \\
\midrule
NumPy (CPU) & CNN cơ bản, 12 epoch & 0,542 & 0,351 & 0,523 & 1,20 tr & 260 s \\
Keras (CPU) & CNN cơ bản & 0,505 & 0,370 & 0,468 & 1,20 tr & 22 s \\
PyTorch (GPU) & CNN cơ bản & \textbf{0,555} & \textbf{0,379} & \textbf{0,589} & 1,20 tr & 12 s \\
\midrule
Keras (CPU) & CNN cải tiến (khối dư) & 0,520 & 0,247 & 0,524 & 1,30 tr & 178 s \\
PyTorch (GPU) & CNN cải tiến (khối dư) & 0,550 & 0,327 & 0,500 & 1,29 tr & 80 s \\
\midrule
\multicolumn{2}{@{}l}{\emph{Tham chiếu (mục trước): ResNet-18 chuyển giao}} & 0,578 & 0,482 & 0,700 & 11,2 tr & 34 s \\
\bottomrule
\end{tabular}
\end{table}

\textbf{Nhận xét.} (i) Ba cách viết cho kết quả cùng một vùng (QWK 0,47–0,59), chênh lệch nằm trong độ dao động của một tập test 400 ảnh với lớp hiếm. Bản NumPy tự viết đạt QWK 0,52, chứng tỏ im2col, col2im và lan truyền ngược đúng, nhưng chậm hơn PyTorch GPU 20 lần. (ii) \emph{CNN cải tiến không tốt hơn} CNN cơ bản trên bộ này, trái với trên MNIST và Fashion-MNIST (\tabref{tab:ch3-domains}). Với 1.080 ảnh huấn luyện, mạng sâu hơn quá khớp nhanh: bản Keras đạt validation tốt nhất ở epoch 3 rồi dừng sớm. (iii) Cả năm mô hình học từ đầu đều kém xa ResNet-18 học chuyển giao (QWK 0,70). Điều này củng cố kết luận của mục trước: với ảnh y khoa ít, \emph{tri thức thị giác tiền huấn luyện} quan trọng hơn cách cài đặt hay độ sâu.

\section{CNN dưới dạng hợp các hàm}
\label{sec:ch3-compose}

CNN cơ bản ở \tabref{tab:ch3-three} là hợp của năm khối hàm. Với ảnh $\mathbf X\in\mathbb R^{48\times48\times3}$:
\begin{equation}
\hat{\mathbf p} = \big(\mathrm{softmax}\circ L_2 \circ \mathrm{ReLU}\circ L_1 \circ \mathrm{vec} \circ P\circ\mathrm{ReLU}\circ C_2 \circ P\circ\mathrm{ReLU}\circ C_1\big)(\mathbf X).
\label{eq:ch3-compose}
\end{equation}
\begin{itemize}
  \item $C_1:\mathbb R^{48\times48\times3}\to\mathbb R^{48\times48\times32}$ và $C_2:\mathbb R^{24\times24\times32}\to\mathbb R^{24\times24\times64}$ là tích chập 3×3: $C(\mathbf X)_{i,j,f}=\sum_{c,u,v}W_{f,c,u,v}\,X_{i+u,j+v,c}+b_f$. Hai tính chất \emph{kết nối cục bộ} và \emph{chia sẻ trọng số} khiến $C_1$ chỉ có $3\cdot3\cdot3\cdot32+32=896$ tham số, trong khi lớp kết nối đầy đủ cùng kích thước cần hơn 200 triệu.
  \item $P$ (MaxPool 2×2) giảm một nửa mỗi chiều không gian và tạo bất biến với dịch chuyển nhỏ: một vi phình mạch dịch đi 1 điểm ảnh vẫn cho cùng đáp ứng.
  \item $\mathrm{vec}$ duỗi khối $12\times12\times64$ thành vector 9.216 chiều; $L_1$ (9.216→128) và $L_2$ (128→5) là các hàm affine, trong đó $L_1$ chiếm 98\% tham số. Đây cũng là chỗ quá khớp, và là lý do CNN cải tiến thay $\mathrm{vec}$ bằng Global Average Pooling.
  \item $\mathrm{softmax}$ cho xác suất 5 mức bệnh.
\end{itemize}
Thứ tự hợp phản ánh \emph{thứ bậc đặc trưng}: $C_1$ học cạnh và đốm màu (mạch máu, xuất huyết điểm), $C_2$ học tổ hợp của chúng (cụm xuất tiết, vùng thiếu máu). Grad-CAM ở mục trước chính là việc lấy gradient của $\hat p_k$ theo đầu ra của khối tích chập cuối, tức là ``nhìn ngược'' qua \eq{eq:ch3-compose}.

\section{Đối chiếu trên ảnh tự nhiên}
\label{sec:ch3-domains}

Cùng hai kiến trúc được huấn luyện trên hai bộ ảnh chuẩn (mã nguồn kế thừa của nhóm): \textbf{MNIST} (70.000 ảnh chữ số viết tay 28×28) và \textbf{Fashion-MNIST} (70.000 ảnh trang phục 28×28, 10 lớp).

\begin{table}[H]
\centering
\caption{Đối chiếu: cùng CNN trên ảnh chữ số và trang phục (10.000 ảnh test mỗi bộ)}
\label{tab:ch3-domains}
\small
\begin{tabular}{@{}llccrr@{}}
\toprule
\textbf{Bộ dữ liệu} & \textbf{Mô hình} & \textbf{Accuracy} & \textbf{Macro-F1} & \textbf{Tham số} & \textbf{Huấn luyện} \\
\midrule
MNIST & CNN cơ bản – NumPy & 0,988 & 0,988 & 421.642 & 1.634 s \\
MNIST & CNN cơ bản – Keras & 0,990 & 0,989 & 421.642 & 97 s \\
MNIST & CNN cơ bản – PyTorch & 0,992 & 0,992 & 421.642 & 22 s \\
MNIST & CNN cải tiến – PyTorch & \textbf{0,997} & \textbf{0,997} & 342.378 & 151 s \\
\midrule
Fashion-MNIST & CNN cơ bản – PyTorch & 0,918 & 0,917 & 421.642 & 26 s \\
Fashion-MNIST & CNN cải tiến – PyTorch & \textbf{0,939} & \textbf{0,939} & 342.378 & 258 s \\
\bottomrule
\end{tabular}
\end{table}

\textbf{Nhận xét.} Trên ảnh tự nhiên với hàng chục nghìn mẫu, mọi kết luận của bài toán đáy mắt đảo ngược: cùng kiến trúc cơ bản đạt 99\% trên MNIST, và CNN cải tiến (khối dư, BatchNorm, tăng cường dữ liệu) hơn bản cơ bản 0,5–2 điểm phần trăm thay vì kém đi. Khác biệt không nằm ở mô hình mà ở \emph{dữ liệu}: ảnh đáy mắt ít mẫu, lớp lệch nặng, và tín hiệu bệnh nằm ở chi tiết nhỏ bị mất khi thu ảnh về 48×48. Vì vậy với ảnh y khoa, học chuyển giao và ảnh độ phân giải cao là hướng cải thiện đúng, chứ không phải chồng thêm lớp.

\section{Thực nghiệm trên Notebook và Demo Localhost}
\label{sec:ch3-notebook-demo}

Notebook \texttt{notebooks/ch3\_cnn\_day\_mat.ipynb} (các ô \texttt{CNN-00} đến \texttt{CNN-04}) cho phép kiểm tra trực tiếp toán học bên dưới lớp tích chập và phân bố tham số của mạng:
\begin{itemize}
  \item \textbf{Ô \texttt{CNN-02} (Kiểm chứng đạo hàm giải tích \texttt{conv\_backward} bằng sai phân hữu hạn)}: Trước khi huấn luyện CNN viết tay bằng NumPy trên 1.080 ảnh đáy mắt, ô \texttt{CNN-02} sinh ngẫu nhiên một lô ảnh nhỏ $X\in\mathbb R^{2\times 3\times 6\times 6}$ và bộ lọc $W\in\mathbb R^{4\times 3\times 3\times 3}$, tính lan truyền tiến bằng \texttt{im2col} và đạo hàm ngược bằng \texttt{col2im}, rồi đối chiếu với đạo hàm số trung tâm $\frac{L(W_{ij}+\varepsilon)-L(W_{ij}-\varepsilon)}{2\varepsilon}$ ($\varepsilon=10^{-5}$). Sai số tương đối cực đại trên toàn bộ trọng số đạt $8{,}28\times 10^{-11}$, chứng minh phép biến đổi hai chiều \texttt{im2col}/\texttt{col2im} bảo toàn chính xác gradient ở biên đệm (\texttt{padding=1}).
  \item \textbf{Ô \texttt{CNN-04} (Giải phẫu hàm hợp 5 khối và nút thắt cổ chai tham số)}: Ô \texttt{CNN-04} tách \texttt{BasicCNN} thành khối trích xuất đặc trưng tích chập (\texttt{features}) và khối phân loại kết nối đầy đủ (\texttt{classifier}). Kết quả in ra cho thấy hai lớp tích chập $3\to 32\to 64$ chỉ dùng đúng $19.392$ tham số ($1{,}62\%$ tổng mô hình), trong khi hai lớp tuyến tính sau \texttt{Flatten} chiếm tới $1.180.421$ tham số ($98{,}38\%$). Con số này giải thích trực tiếp vì sao các kiến trúc hiện đại (như CNN cải tiến và ResNet-18) thay \texttt{Flatten} bằng Global Average Pooling để giảm hơn $90\%$ tham số ở đầu phân loại và hạn chế quá khớp khi dữ liệu y khoa nhỏ.
\end{itemize}

Trên ứng dụng web cục bộ tại \texttt{http://localhost:8095} (\figref{fig:app_ch3_retina}), tab \textbf{Ch.3 — CNN · Võng mạc ĐTĐ (RetinaMNIST)} nạp đồng thời cả 5 mô hình đã huấn luyện (\texttt{scratch}, \texttt{tf\_basic}, \texttt{torch\_basic}, \texttt{tf\_improved}, \texttt{torch\_improved}), cho phép chọn mẫu ảnh đáy mắt kiểm tra và so sánh trực quan phân phối xác suất 5 mức độ bệnh võng mạc đái tháo đường cùng thời gian suy luận của từng bản cài đặt.

\figw[0.92\linewidth]{app_ch3_retina}{Giao diện thử nghiệm Chương 3 trên \texttt{http://localhost:8095}: so sánh dự đoán 5 mức bệnh võng mạc đái tháo đường của CNN NumPy, Keras và PyTorch trên cùng ảnh đáy mắt RetinaMNIST}

